\section{Test-Plan per la funzionalità di aggiunta commento}
In questa sezione viene descritto il piano di test adottato. Il piano di test include le 
strategie, le metodologie e gli strumenti utilizzati per 
garantire la qualità del software sviluppato.
In particolare andremo a descrivere il test plan per la funzionalità di aggiunta di un commento
nella sezione dei dettagli di una Issue.

\subsection{Strategia di Test}
La strategia di test adottata per la funzionalità di aggiunta commento prevede l'utilizzo di test 
automatizzati. I test automatizzati verranno eseguiti utilizzando jUnit per verificare che la funzionalità
di aggiunta commento funzioni correttamente in diversi scenari. In particolare la strategia utilizzata è di Black-Box 
testing, questa si basa sul partizionamento in classi di equivalenza, concentrandosi sui requisiti funzionali
senza considerare l'implementazione interna del codice. Per garantire una copertura soddisfacente, abbiamo utilizzato
l'approccio di testing R-WECT (Robust-Weak Equivalence Class Testing), che prevede un'analisi approfondita delle classi non valide, 
e una più superficiale delle classi valide per mantenere un numero gestibile di test case.

\subsection{Strumenti utilizzati}
Per l'esecuzione dei test automatizzati, utilizziamo jUnit, un framework di testing per Java che consente di scrivere e
eseguire test in modo semplice ed efficiente. jUnit offre funzionalità avanzate per la gestione dei test, 
inclusa la possibilità di raggruppare i test in suite e di generare report dettagliati sui risultati dei test.
Altro strumento utilizzato è Mockito, un framework di mocking per Java che consente di creare oggetti fittizi (mock) per isolare le unità di codice durante i test. 
Mockito permette di simulare il comportamento di dipendenze esterne, facilitando la scrittura di test unitari più efficaci.

\newpage

\section{Test Case}
Di seguito sono riportati i test case per la funzione di aggiunta commento (metodo addComment) e la funzione di validazione dei dati
utente (metodo validateUserData).

\subsection{Test Case per il metodo addComment}
\subsubsection{Classi d'equivalenza individuate}
Queste sono le classi di equivalenza individuate per il metodo \texttt{addComment}:
\begin{itemize}
    \item \textbf{Oggetto DTO}:
    \begin{itemize}
        \item \textbf{CE1 (Valida)}: \texttt{DTO} istanzato.
        \item \textbf{CE2 (Non Valida)}: \texttt{DTO} nullo (\texttt{null}).
    \end{itemize}
    \item \textbf{Campo \texttt{text}}:
    \begin{itemize}
        \item \textbf{CE3 (Valida)}: Stringa valida (non vuota).
        \item \textbf{CE4 (Non Valida)}: \texttt{null}.
        \item \textbf{CE5 (Non Valida)}: Stringa vuota (\texttt{""}).
        \item \textbf{CE6 (Non Valida)}: Stringa blank (\texttt{" "}).
    \end{itemize}
    \item \textbf{Campo \texttt{issueId}}:
    \begin{itemize}
        \item \textbf{CE7 (Valida)}: Intero positivo ($> 0$).
        \item \textbf{CE8 (Non Valida)}: \texttt{null}.
        \item \textbf{CE9 (Non Valida)}: Zero ($0$).
        \item \textbf{CE10 (Non Valida)}: Negativo ($< 0$).
    \end{itemize}
    \item \textbf{User (DB/Dao)}:
    \begin{itemize}
        \item \textbf{CE11 (Valida)}: Utente esistente (trovato).
        \item \textbf{CE12 (Non Valida)}: Utente inesistente (non trovato/\texttt{null}).
    \end{itemize}
\end{itemize}

\subsection{Test Case per il metodo validateUserData}
\subsubsection{Classi d'equivalenza individuate}
Queste sono le classi di equivalenza individuate per il metodo \texttt{validateUserData}:
\begin{itemize}
    \item \textbf{Tutti i campi String}:
    \begin{itemize}
        \item \textbf{CE1 (Valida)}: Stringa valida (non \texttt{null}, non vuota).
        \item \textbf{CE2 (Non Valida)}: Valore \texttt{null}.
        \item \textbf{CE3 (Non Valida)}: Stringa vuota (\texttt{""}).
        \item \textbf{CE4 (Non Valida)}: Stringa solo spazi (\texttt{" "}).
    \end{itemize}
    \item \textbf{Role}:
    \begin{itemize}
        \item \textbf{CE5 (Valida)}: Ruolo "Admin".
        \item \textbf{CE6 (Valida)}: Ruolo "Normal".
        \item \textbf{CE7 (Valida)}: Ruolo "External".
        \item \textbf{CE8 (Non Valida)}: Ruolo non riconosciuto.
        \item \textbf{CE9 (Non Valida)}: Ruolo \texttt{null}.
    \end{itemize}
\end{itemize}

\newpage
  


